\section{Multi-Agent Fine-Tuning: prolongar el autoentrenamiento con reasoning diverso}

\subsection{Por qué el synthetic data de un solo modelo se estanca}

\videofigure{figures/multiagent-finetuning.jpg}{Multi-Agent Fine-Tuning usa varios generation agent y critic agent para producir, discutir y filtrar synthetic reasoning data.}{00:07:28--00:08:42}

\videofigure{figures/diversity-bottleneck.jpg}{Cuando un solo modelo genera repetidamente sus propios datos de entrenamiento, el reasoning se vuelve cada vez más homogéneo y las ganancias por iteración se estancan pronto.}{00:08:42--00:09:58}

Los datos de pretraining son generados por una gran cantidad de humanos a lo largo del tiempo, y cubren de manera natural distintos estilos, conocimientos y rutas de pensamiento; un solo LLM, incluso al subir la temperature, suele muestrear alrededor de unas pocas plantillas de alta probabilidad. El rejection sampling solo conserva las respuestas correctas, lo cual comprime aún más la distribución y produce mode collapse.

Se puede denotar la distribución de synthetic data de la ronda $m$ como $q_m(r\mid x)$; si el modelo siempre se hace fine-tuning sobre sus propias salidas:

$$
q_{m+1}(r\mid x)\propto q_m(r\mid x)\,w(r),
$$

donde $w(r)$ es un peso de filtrado por corrección o preferencia. Si $w$ siempre favorece los patrones de alta probabilidad ya existentes, el reasoning de baja probabilidad pero válido irá desapareciendo ronda tras ronda.

\begin{warningbox}{Temperature no es lo mismo que diversidad semántica}
Una temperature alta puede aumentar la variación de la redacción y de los tokens locales, pero no necesariamente produce estrategias de solución distintas. La diversity realmente útil debe medirse en el nivel de las hipótesis, la descomposición, los algoritmos y las rutas de evidencia.
\end{warningbox}

\subsection{Generation agents y critic agents}

\videofigure{figures/generate-data.jpg}{Varios generation agent responden primero de manera independiente, luego reúnen las respuestas de los demás agents y entran a la siguiente ronda de discusión.}{00:09:58--00:10:56}

\videofigure{figures/finetune-generators.jpg}{El generation model continúa el fine-tuning sobre pares prompt--response de alta calidad que coinciden con el resultado mayoritario.}{00:10:56--00:11:26}

\videofigure{figures/finetune-critics.jpg}{El critic model aprende a distinguir las trayectorias que fueron correctas desde el inicio de aquellas que se corrigieron solo tras la discusión.}{00:11:26--00:11:58}

Supóngase que hay $N$ generators, y que la respuesta de la ronda $t$ es $a_i^{(t)}$. El sistema incorpora a la siguiente ronda el resumen de las respuestas de los demás agents $s_{-i}^{(t)}$ y la retroalimentación del critic $c_i^{(t)}$:

$$
a_i^{(t+1)}\sim p_{\theta_i}
\left(a\mid x,a_i^{(t)},s_{-i}^{(t)},c_i^{(t)}\right).
$$

\begin{itemize}
    \item $x$: la pregunta original;
    \item $a_i^{(t)}$: la respuesta actual del $i$-ésimo generator;
    \item $s_{-i}^{(t)}$: el resumen de los puntos de vista de los demás agents;
    \item $c_i^{(t)}$: la retroalimentación del critic sobre las deficiencias de la respuesta.
\end{itemize}

\videofigure{figures/multiagent-debate.jpg}{Varias rondas de debate hacen que los specialized agents expongan mutuamente sus puntos ciegos, y al final se usa majority vote para producir las etiquetas de entrenamiento.}{00:11:58--00:13:08}

\begin{importantbox}{El beneficio de tener varios agents proviene de un inductive bias diferenciado}
Si todos los agents son el mismo checkpoint, con el mismo prompt y la misma configuración de decoding, compartirán los mismos modos de falla. Es necesario generar una especialización estable mediante subconjuntos de datos, roles, hiperparámetros o familias de modelos distintos.
\end{importantbox}

\subsection{Diversidad y resultados de mejora continua}

\videofigure{figures/diversity-results.jpg}{La embedding dissimilarity muestra que el esquema multi-agent logra mantener una mayor diversidad de respuestas incluso después de varias rondas de fine-tuning.}{00:13:08--00:13:38}

\videofigure{figures/multiagent-results.jpg}{En varios modelos de código abierto, el multi-agent fine-tuning sigue aumentando la accuracy con las iteraciones, mientras que el esquema single-agent se estanca o decae antes.}{00:13:38--00:14:08}

\videofigure{figures/generalization-results.jpg}{El reasoning generado por varios agents no solo mejora el dominio de entrenamiento, sino que también muestra transferencia en benchmarks matemáticos relacionados.}{00:14:08--00:15:02}

\begin{knowledgebox}{Es necesario monitorear al mismo tiempo el performance y la diversity}
Si la accuracy sube pero la embedding diversity cae rápidamente, esto puede significar que el sistema está agotando su espacio de exploración futuro; si la diversity es muy alta pero la accuracy no mejora, puede tratarse solo de ruido sin utilidad. Ambas deben conformar una Pareto frontier.
\end{knowledgebox}

\subsection{Resumen de la sección}

Multi-Agent Fine-Tuning usa specialized generators para ampliar la distribución del reasoning, y usa critics junto con votación mayoritaria para controlar la calidad, mitigando así la homogeneización y el estancamiento por iteración del synthetic data de un solo modelo.
